\scanpage{34}

\uline{Remarque 1} La \uline{propriété $LG$ faible} \uline{équivaut} à $(LG)_0$.

En effet, il suffit de prouver qu'elle implique que $H^{\unsure{1}}(X,F)\to H^1(X_{/Y},F_{/Y})$ est injectif, i.e. que tt ext. de $\underline{\mathcal O}_X$ par $F$ qui splitte formellement, splitte régulièrement, ce qui est immédiat. [Question : Si $\unsure{\alpha|\ast}$ est un isom pour $i\leq n$, tt $F$ loc libre, est-il vrai que $\unsure{\alpha|\ast}$ est un mono pour \deleted{\unsure{tout}} $i=n+1$, tt $F$ loc libre ?.] Oui, prendre \inserted{une \deleted{\unsure{résolution}} \inserted{\unsure{résol. univers.|résolution}} \unsure{singulière|injective|régulière|injection}, $0\to F\to I\to G\to0$, et appliquer le foncteur $_{/Y}$, et le lemme des cinq]}

\uline{Remarque 2} Soit $(L_\alpha)$ une famille de Modules loc. libres sur $X$, telle que tt module loc. libre sur $X$ soit \inserted{$\simeq$} quotient d'une somme directe de Modules $\check L_\alpha$, i.e. $\simeq$ au \unsure{ann. dual|ann. d'un|ss. mod.|sous-module} \inserted{loc. facteur direct} d'une somme directe de Modules $L_\alpha$ $[$Ex : $1^\circ$ $X$ quasi \deleted{projectif} \inserted{affine}, $(L_\alpha)$ réduit : $\underline{\mathcal O}_X$ $2^\circ$ $X$ quasi-projectif, $(L_\alpha)$ formé des $\underline{\mathcal O}(m)$, $m\geq m_0$ $]$.

Alors $(LG)_n$ équivaut à la condition

\[
H^i(X,L_\alpha)\overset{\sim}{\longrightarrow}H^i(X_{/Y},L_{\alpha/Y})
\quad\left|\begin{aligned}&\text{isom pour }i\leq n\\[-2pt]&\text{inj pour }i=n+1\end{aligned}\right.
\]
% slanted
\margin{N.B. Le même argument marche si on n'exige pas la condition $H^{n+1}(\unsure{F|X})\to H^{n+1}(\hat{\unsure{F|X}})$ \uline{injectif}.}

\deleted{\unsure{Cet énoncé|Aut. rédac}} \deleted{[N.B. \unsure{Conclusion :|Contrairement :} \unsure{pour la|pourquoi} preuve \unsure{on mettra} une condition sur $i=n+1$ !]}. Dém par récurrence sur $n$, trivial si $n<-1$, ds $n+1<0$, supposons \deleted{\unsure{un}} que c'est prouvé

% archive sheet number: 33
